% ---------------------------------------------------------------------------
% preface.tex
% ---------------------------------------------------------------------------
\chapter*{Preface}
\addcontentsline{toc}{chapter}{Preface}
\markboth{Preface}{Preface}

Reading rules is not the same as reading code. This booklet is one design,
worked all the way through: a CORDIC rotator, in VHDL and in SystemVerilog,
verified by six testbenches in three languages. Every file is reproduced in
full, not excerpted.

This is the \textbf{Example Cookbook}, one of four companions covering
\tool{SystemVerilog for VHDL Designers}. It assumes the material of
\refdesignbook, \refverificationbook and \refsimulationbook, and applies all
three, in one place, to one design.

\section*{What this booklet is}

\Cref{ch:case} develops \code{cordic\_rot} from its specification and
fixed-point format, through a VHDL implementation and a SystemVerilog one,
to six testbenches: a plain VHDL testbench, a class-based SystemVerilog
testbench, a UVM environment, a C++ testbench driven through
\tool{Verilator}, and two Python testbenches, one talking to the design over
a socket and one written with \tool{cocotb}. Every construct used has
already been introduced somewhere in the companion booklets; here you see it
used together with everything else, under time pressure, in a design with
real corners: a state machine, a parameterised datapath, a fixed-point
number format, and a handshake protocol.

\section*{The code}

All the sources are distributed with this document under \file{code/cordic/}.
Every one of them has been analysed and simulated before being typeset, and
every line of tool output quoted in this booklet was produced by actually
running something. Where a construct does not work in a particular tool, the
text says which tool and which version.

\section*{Acknowledgement}

The choice of a CORDIC rotator as the worked example is not accidental. It
is small enough to read in one sitting, large enough to need a state
machine, a parameterised datapath, a fixed-point format and a handshake
protocol, and arithmetic enough that a golden model in Python earns its
keep. If you understand every line of this booklet, you can write
SystemVerilog.
